\section{Overview}

\subsection*{Proposed thesis structure}
\begin{enumerate}
    \item \textbf{Introduction/Theory} 
    \begin{itemize}
        \item Ion trap theory needed. 
        \begin{itemize}
            \item Trapping fields
            \item QP transitions + group theory on this
            \item Ion-light interactions. With spin, with motion.
            \item MS gates + SDFs
        \end{itemize}
        \item Noise (Look at Harty + Amy thesis)
        \begin{itemize}
            \item Model for general noise. Using error channels + PTM + Choi matrices. Discuss classifications of noise channels.
            \item How to measure noise. Tomography, RB.
            \item How to simplify noise using RB.
            \item How k-copy RB simplifies noise but retains structure.
            \item Group + representation theory for RB
            \item Reducing physical noise. Coherent with H engineering, incoherent with fast gates.
            \item Reducing the impact of noise by error mitigation techniques.
            \item Virtual distillation on spatially correlated noise.
        \end{itemize}
    \end{itemize}

    \item \textbf{Experimental Apparatus} 
    -- Description of the new experimental apparatus. 
    \begin{itemize}
        \item Block diagram of whole experiment highlighting automation.
        \item See current chapter for most sections.
        \item Ground loops and how we found them. Suggest improvements to current system.
        \item Single addressing system. Chat to IRO about what minimal info to include.
        \item Section on Pulse shaping with phaser, sinara?
        \item Verdi noise spectra
        \item 729 laser characterisation. Including linewidth, cavity ringdown, noise spectra. FNC noise reduction. Keeping it alive? Autolock feature.
    \end{itemize}

    \item \textbf{Experimental Characterisation} 
    -- spin/motion/spin-motion elementary interactions that are used throughout the thesis. This is the chapter submitted for confirmation of status.\\
    \begin{itemize}
        \item Spin coherence times with and without box, and with and without echo sequence.
        \item 729 laser characterisation. Including linewidth, cavity ringdown, noise spectra. FNC noise reduction. Keeping it alive?
        \item New mode stability with enclosed resonator
        \item Motional coherence time
        \item Going to long chains
        \item micromotion compensation
        \item All mode frequencies. Using James.
        \item All modes SBC. Here also a discussion on sub-doppler cooling options.
        \item MS gates improved + Walshing
        \item Error budget for MS gates
        \item QP matrix elements -> polarisations -> Stark shifts?
    \end{itemize}

    \item \textbf{k-copy Randomised Benchmarking} 
    -- Presentation of the theory and experimental results demonstrating
    parallel single-qubit randomised benchmarking to distinguish local and
    global noise contributions on each qubit. This is the paper extended.
    \begin{itemize}
        \item How k-copy RB simplifies noise but retains structure.
        \item Group + representation theory for k-copy RB to derive correlated populations
        \item Shared global fields leading to spatially correlated noise. Error channel for this.
        \item Optimal sampling for multi-exponential
        \item Parameter range
        \item Noise injection calibration
        \item Stark shift calibration
        \item Extracting errors for large parameter space
        \item Parity measurements
        \item Pairwise measurements
    \end{itemize}

    \item \textbf{Virtual Distillation}
    -- Presentation of the experimental results demonstrating virtual distillation.
    \begin{itemize}
        \item two copy, B matrix demonstration
        \item any-to-any two-qubit gate
        \item pi/4 calibration
        \item 
    \end{itemize}

    \item \textbf{Conclusion and Outlook} 

\end{enumerate}